\section{Introduction}
Young adults from better-off families are much more likely to attend college than those from worse-off families. For example, in the United States children from families where at least one parent has attained higher education are 37 percentage points more likely to have a college degree than children from families where neither has; the gap is similarly high in other industrialized economies (\citealp{oecd2024eagdata}). One policy response to this intergenerational inequality is to provide college admission advantages to students from disadvantaged contexts. Context-based admissions are gaining increasing attention, especially as admissions based on race or ethnicity are proving contentious and have been severely restricted in some countries (\citealp{arcidiacono2016affirmative,feingold2023affirmative}). \par 



By altering the link between academic effort and admission chances, preferential admission policies change the study incentives of disadvantaged students while still in school (\citealp{coate1993will}). This paper analyzes how students’ subjective beliefs about their admission chances and their future college success shape these incentive effects, and how study effort responses, in turn, influence who enters college under preferential admissions and how academically prepared they are. Understanding how students' perceptions shape the impacts of admission policies is essential for designing them effectively.\par 





We study these questions in the context of Chile, which is uniquely well-suited for four reasons. First, it introduced a nationwide policy, PACE ({\itshape Programa de Acompañamiento y Acceso Efectivo a la Educación Superior}), that granted large admission advantages to students from disadvantaged schools. Second, Chile operates a transparent, centralized college admission system, allowing us to observe the actual incentives students face. Third, the country maintains rich longitudinal administrative data that track students from high school through college enrollment, which we could link to survey data on students’ beliefs. Fourth, the rollout of PACE was randomized across high schools, allowing us to identify the policy's impacts.\footnote{One of the paper’s authors, Michela Tincani, led the experimental evaluation of PACE in collaboration with Chile’s Ministry of Education and Ministry of Finance (\citealp{dipres2022experimental}). Policy reports on the experimental evaluation of the program include \citealp*{evaluacion_mineduc_2019,evaluacion_mineduc_2020,evaluacion_mineduc_2022}.}\par 




PACE targets students in disadvantaged high schools and guarantees admission to those graduating in the top 15\% of their class to colleges participating in the centralized system, waiving the national entrance exam requirement. These colleges offer five-year (or longer) academically oriented programs, and their commitment to the policy is formalized through agreements with the government. PACE does not replace the regular route: students in PACE schools may still sit the national exam and compete for seats through the regular channel; the policy offers an additional pathway for top-ranked students. Students attending PACE high schools are substantially more disadvantaged than typical college entrants: their $10^{th}$-grade standardized test scores are, on average, 1.4 standard deviations lower, and their household income is roughly a third as high. PACE thus expanded access to selective colleges for a population dramatically underrepresented in the college system.\par 





We constructed a new longitudinal dataset that links high-quality administrative records with original survey data collected in schools. The data follow 9,006 students who were in 11\textsuperscript{th} grade in 2016 across 128 high schools—half randomly assigned to receive the PACE program and half to serve as controls. The administrative records cover students from 9\textsuperscript{th} grade through five years after high school, and include detailed measures of academic performance, grades, demographics, and higher education outcomes, including enrollment and persistence.\par 

To complement these administrative data, we designed and administered surveys to students in their final year of high school. The surveys capture students’ beliefs about their academic ability (both absolute and relative), their perceived returns to effort, their expectations about college performance, and the monetary returns to college. We also collected data on self-reported effort and administered a standardized test to measure academic achievement. To understand how schools may respond to preferential admissions, we also surveyed teachers and principals about instructional focus, grading practices, and support classes. We linked all survey responses to the administrative data via unique student, classroom, and school identifiers.\par 


There are two main experimental findings. First, PACE increased college admissions and enrollments among disadvantaged students by 4.1 and 3.0 percentage points— 36\% and 35\% increases relative to the control group. These effects were concentrated among students who, in 10\textsuperscript{th} grade (before the experiment started), ranked in the top 15\% of their high school cohort; students in the bottom 85\% experienced no significant change in admissions or enrollment. While the initial effects were substantial, they declined over time: five years after high school, the impact on continuous enrollment or graduation from a selective college was 1.5 percentage points—a 30\% increase relative to control students. Second, PACE reduced students’ study effort and achievement in high school by 0.1 standard deviations. In contrast to the enrollment gains, these reductions were widespread, impacting students across the achievement distribution.\footnote{The PACE policy was introduced by the Bachelet government and retained by the subsequent Pi\~{n}era government after reviewing early experimental evidence on college admissions and enrollment impacts (\citet*{evaluacion_mineduc_2019}).} \par 




To understand the waning enrollment impacts, we examine the type of selective colleges students attend and the characteristics of college entrants. We find no systematic changes in the selectivity, field of study, or geographic location of the programs students attend, suggesting that college match does not mediate the waning enrollment impacts over time. Instead, we cannot rule out changes in the composition of college entrants in terms of unmeasured baseline ability, nor differences in academic preparedness. Supporting the latter, we show that effort and achievement in the last high school year predict college persistence, suggesting that reduced effort and achievement in high school may have left PACE students less prepared for college, contributing to lower persistence over time.\par 

Next, we examine the mechanisms behind the observed reduction in pre-college effort. We find no evidence that PACE changed instructional practices or school-level academic support, nor did it affect students’ perceptions of the monetary returns to college. These results suggest that neither school-side adjustments nor updated beliefs about the value of college explain the observed reduction in effort. Instead, we find evidence consistent with students responding to perceived incentives under the new admissions regime. By linking survey responses to actual academic outcomes in the administrative data, we document systematic over-optimism about both absolute and relative ability—suggesting that many students misperceived how close they were to the regular and preferential admission cutoffs. Effort reductions were concentrated among students who believed they were well above the PACE admissions threshold—consistent with the perception that PACE lowered the returns to effort in securing college admission. Additional belief data show that students were also over-optimistic about their likelihood to persist in college, and did not view high school effort as important for succeeding in college—suggesting they perceived little consequence for their future college success from reducing effort. \par 

Motivated by these findings, we develop a dynamic structural model of students’ educational choices to examine the implications of these belief distortions for the behavior of students and for how they would react to new, alternative PACE designs. In the model, students with different observed and unobserved characteristics make pre-college effort, entrance-exam taking, and enrollment decisions based on subjective beliefs about the returns to pre-college effort in securing regular and PACE admissions and in persisting in selective college. The model also includes objective admission and persistence likelihoods. In the model, belief-driven pre-college decisions can shape long-term college outcomes by affecting both who enters selective college and their likelihood to persist. We exploit our strategically designed survey questions to identify subjective returns to effort, and use experimental variation to help identify objective returns. The estimated model can match both targeted and untargeted treatment‐effect patterns.\par










The first model result is that 77\% of the observed association between pre-college effort and college persistence reflects a causal effect, while the remainder is due to unobserved variable bias—students more likely to persist also tend to exert more effort in school.\par 




Second, we simulate a counterfactual in which both control and treated students have rational expectations to assess the role of belief distortions. Relative to the baseline with biased beliefs, students in both groups would exert less pre-college effort because they would no longer overestimate its returns in securing regular admission nor their likelihood of college persistence. Under rational expectations, PACE would not have reduced pre-college effort, as it would have slightly increased its returns by making college more attainable. Nevertheless, the PACE effect on enrollment would have weakened, as students would have correctly anticipated lower persistence and valued college entry less. Subjective beliefs, therefore, fundamentally shaped the effects of PACE.\par 





Third, we simulate the effects of pairing PACE with an informational intervention that corrects the beliefs of treated students only.
Although belief distortions help explain the decline in pre-college effort under PACE, correcting treated students' overoptimism about their admission and persistence chances would not eliminate this unintended effect—it would amplify it. Students who receive PACE with belief correction exert even less effort than students who only receive PACE. And while correcting beliefs improves the composition of college entrants from PACE schools in terms of baseline test scores, it ultimately reduces their persistence through lower pre-college effort. As a result, the treatment effects of this counterfactual policy are more negative for pre-college effort and smaller for admissions and long-term enrollment than the treatment effects of PACE alone.\par 


A government interested in avoiding PACE's unintended impacts on pre-college effort could instead pair PACE with an intervention informing students about the role of pre-college effort in supporting college success, without correcting their other over-optimistic beliefs. This design mitigates the decline in effort without dampening enrollment gains. This suggests that which misperceptions are addressed can shape the impacts of preferential admissions.\par


Finally, we use the model to simulate the impacts of alternative top-percent cutoffs for preferential admissions. More generous cutoffs lead to larger gains in enrollment and persistence, but they also generate increasing numbers of dropouts. Beyond the current 15\% cutoff, the effect on the number of college dropouts would exceed the effect on the number of college enrollees on track to graduate.\par 


This paper contributes to the literature on preferential college admissions by providing unified evidence on how a preferential admission policy in Chile affected students' outcomes from before college entry to five years post high school. Two separate strands of the literature have documented impacts on pre-college academic outcomes (\citealp{golightly2019does,akhtari2024affirmative,khanna2020does}) and on longer-term college enrollment and persistence (\citealp{long2010policy}, \citealp{niu2010impact}, \citealp{daugherty2014percent}, \citealp{bleemer2021top}, \citealp{black2023winners}).\footnote{See also \citealp{hastings2012effect} for related evidence on K-12 responses to future educational opportunities, and \citealp{arcidiacono2015affirmative} for a broad review of the literature on affirmative action in college admissions.} This paper shows that the incentive effects on pre-college outcomes are not only policy-relevant per se, but they also matter for preferential admissions' longer-term impacts on college persistence. Additionally, the paper adds experimental evidence, which has so far remained limited.\par 

Second, the paper contributes new evidence on how preferential admissions affect students and schools before college. Using linked administrative and large-scale survey data, we examine students, teachers, and school-level inputs. These data show that students’ beliefs are central to how they respond to preferential admissions. Prior work has shown that students’ biased beliefs about their admissions chances affect application decisions (\citealp{larroucau2024college, hakimov2025confidence}); this paper shows that such beliefs also distort human capital investments before college in response to admissions policy.\par 

Third, the paper contributes to the structural literature on preferential admissions (e.g., \citealp{arcidiacono2005affirmative,kapor2024transparency,otero2023affirmative}) by developing and estimating a dynamic model that endogenizes pre-college effort, incorporates subjective beliefs, and leverages experimental variation in estimation. While a small number of recent models incorporate effort decisions (i.e., \citealp{hickman2024pre,grau2018impact,borghesan2022heterogeneous}), none of them have modeled the role of beliefs or used randomized variation in estimation. By relaxing rational expectations, the paper also contributes to a broader structural literature on beliefs in education. Existing work has focused on information frictions during college (e.g., \citealp{stinebrickner2014major,wiswall2015determinants, arcidiacono2020ex}), whereas this paper examines frictions before college and shows how they shape long-run educational outcomes. Related studies have used belief data to estimate static school-choice models (e.g., \citealp{bobba2025perceived, kapor2020heterogeneous}), or incorporated expectations about future choices in dynamic models (e.g., \citealp{van2012use, delavande2019university}). Finally, by combining experimental and structural methods, the paper contributes to the growing literature that uses Randomized Controlled Trials to discipline structural models (e.g., \citealp{todd2006assessing,todd2020best,attanasio2011education}).





